% CS6400-A Project Milestone (Group 18)
% Reuses the ACM sigconf template from the project proposal.
% The `changes` package marks every deviation from the proposal, as the
% milestone instructions request.
%
% TODO before submitting:
%   1. Answer generation currently names qwen3:8b served with Ollama; change
%      Sections 3.1 and 3.3 together if the team picks a different model.
%
% Verified with tectonic 0.17 (XeTeX): 5 pages, no errors, no overfull hbox.

\documentclass[sigconf,nonacm]{acmart}

% commandnameprefix=ifneeded: acmart already defines \comment, which changes
% would otherwise redefine and abort on.
\usepackage[commandnameprefix=ifneeded]{changes}
\usepackage{booktabs}
\usepackage{array}
% Long identifiers and a narrow two-column measure leave TeX no good break;
% let it stretch the spaces rather than overrun the column.
\setlength{\emergencystretch}{3em}
\newcolumntype{L}[1]{>{\raggedright\arraybackslash}p{#1}}

\begin{document}

\title{CS6400-A Project Milestone (Group 18)}
\subtitle{Latency-Aware Adaptive Reranking for Hybrid RAG Retrieval}

\author{Wei-Po Lin}
\email{wlin401@gatech.edu}
\affiliation{\institution{Georgia Institute of Technology}\country{}}

\author{Yu-An Chen}
\email{ychen3606@gatech.edu}
\affiliation{\institution{Georgia Institute of Technology}\country{}}

\author{You-Hsiang Tsai}
\email{ytsai309@gatech.edu}
\affiliation{\institution{Georgia Institute of Technology}\country{}}

\author{Siang-Yu Cheng}
\email{scheng371@gatech.edu}
\affiliation{\institution{Georgia Institute of Technology}\country{}}

\maketitle

\section{Progress Check In}

\subsection{Implementation Progress}

We break the project into nine components. Five are complete and four have not
started.

\begin{table}[h]
\centering
\small
\begin{tabular}{@{}L{2.5cm}L{1.5cm}L{3.4cm}@{}}
\toprule
\textbf{Component} & \textbf{Status} & \textbf{Notes} \\
\midrule
Dataset processing      & Complete    & Parsing, chunking, metadata, provenance IDs \\
Evidence mapping        & Complete    & Gold evidence text mapped to paragraph and chunk IDs \\
Sparse retrieval (BM25) & Complete    & Original implementation, standard library only \\
Dense retrieval         & Complete    & \texttt{BAAI/bge-m3} embeddings, FAISS exact index \\
Retrieval evaluation    & Complete    & Recall@K, MRR, hit rate, latency; reproducibility manifests \\
RRF hybrid fusion       & Not started & Next component \\
Adaptive reranking      & Not started & Cross-encoder plus confidence gating \\
Secondary benchmark     & Not started & BEIR-SciFact preparation \\
Answer generation       & Not started & QA accuracy is a secondary metric \\
\bottomrule
\end{tabular}
\end{table}

\paragraph{Dataset processing} We prepare QASPER v0.3 into four artifacts: source
paragraphs, retrieval chunks, questions, and reference annotations. Chunks are
180 words with 30 words of overlap and never cross a paragraph boundary, so every
chunk carries the identifier of the paragraph it came from. The train split
yields 888 papers, 47{,}770 paragraphs, 50{,}189 chunks, and 2{,}593 questions.
Every build records the input checksum, the chunking parameters, the counts, and
the output checksums in a manifest, and the scripts refuse to overwrite existing
outputs. This works as expected: an independent rebuild of the train split on a
second machine reproduced the recorded corpus checksum exactly.

\paragraph{Evidence mapping} QASPER annotates gold evidence as paragraph text
rather than identifiers, so a separate step matches each evidence string back to
a prepared paragraph. Matches are classified as exact, whitespace-normalized,
ambiguous, figure or table, section heading, unmatched, or empty. Only
unambiguous matches on answerable annotations become retrieval positives. The
step also emits an audit file of every entry that needs attention, which is what
our eligibility policy in Section~\ref{sec:setup} is built from. Working as
expected.

\paragraph{Sparse retrieval} BM25 with $k_1 = 1.5$ and $b = 0.75$ over a
Unicode word tokenizer, implemented from scratch rather than through a library,
with deterministic tie-breaking by chunk identifier. Index construction takes
1.9~seconds over 50{,}189 chunks and 47{,}344 distinct terms; retrieving the top
100 chunks for all 2{,}593 questions takes 26.4~seconds. Working as expected.

\paragraph{Dense retrieval} \texttt{BAAI/bge-m3} embeddings in a FAISS
\texttt{IndexFlatIP} index over normalized vectors, so scores are exact cosine
similarities with no approximation. The builder refuses to run if any chunk would
exceed the encoder's token limit; the longest QASPER chunk is 1{,}061 tokens
against a limit of 8{,}192, so nothing is truncated. Working as expected.

\paragraph{Retrieval evaluation} A single evaluator scores any number of run
files in one invocation, which forces every compared system to share one eligible
question set, one chunk-to-paragraph mapping, and one latency estimator. It
verifies the full provenance chain by checksum --- corpus to index to run --- and
refuses to score runs produced from a different query set. Three regression tests
cover the cases that would silently corrupt the results table. Working as
expected.

\subsection{Division of Work}

\begin{table}[h]
\centering
\small
\begin{tabular}{@{}L{2.3cm}L{2.0cm}L{3.1cm}@{}}
\toprule
\textbf{Member} & \textbf{Role} & \textbf{Completed work} \\
\midrule
1. Yu-An Chen & Data and Preprocessing & Repository setup, QASPER parsing and chunking, evidence mapping, preprocessing documentation, exploratory data analysis \\
2. Wei-Po Lin & Retrieval and Indexing & BM25 index construction and retrieval, sparse pipeline tests \\
3. Siang-Yu Cheng & Reranking and Adaptive Policy & Dense embedding generation, FAISS index construction, dense retrieval \\
4. You-Hsiang Tsai & Evaluation and Generation & Evaluation harness, retrieval metrics, latency measurement, baseline experiment runs, result analysis, reproducibility manifests \\
\bottomrule
\end{tabular}
\end{table}

The roles are those assigned in the proposal, with one scheduling adjustment.
Both required baselines were built in parallel: Member~2 implemented sparse
retrieval while Member~3 implemented dense retrieval, rather than Member~2
building both in sequence. This brought the milestone's hard requirement forward
and let evaluation begin earlier. Member~3 now moves to the reranking and
adaptive policy that the proposal assigns to that role, and Members~2 and~3
jointly own RRF fusion, which sits between their two components.

\subsection{Progress Assessment and Remaining Work}
\label{sec:plan}

Measured against what our proposal committed to deliver by the milestone --- a
working preprocessing pipeline, both required baselines, and preliminary
retrieval results --- we are on track. All three exist and are evaluated.

Measured against the proposal's internal schedule we are behind on two items.
RRF fusion was scheduled for Milestone~2 (Oct.~13--26) and is not implemented.
Preparation of the secondary benchmark was scheduled for Milestone~1
(Oct.~1--12) and has not begun. Both slipped for the same reason: establishing a
defensible evaluation methodology took longer than planned. QASPER annotates
evidence as free text at paragraph granularity while retrieval operates on
overlapping chunks, and resolving that mismatch correctly --- rather than
reporting a number that silently double-counts overlapping chunks --- consumed
the time budgeted for fusion. We consider this a worthwhile trade, because every
number in the rest of this project depends on that decision being right.

The remaining work, with the deadline held fixed at Dec.~15:

\begin{itemize}
\item \textbf{RRF fusion} (Nov.~1--5). Fuse the existing BM25 and dense run
files. Both are already written in an identical schema over an identical corpus
and query set, so fusion consumes run files rather than re-running retrieval.
\item \textbf{Confidence signals} (Nov.~5--9). Retriever agreement, measured as
overlap between the two top-$K$ lists, and ranking confidence, measured as score
separation among the top fused candidates.
\item \textbf{Cross-encoder integration and always-rerank baseline}
(Nov.~9--16).
\item \textbf{Adaptive policy and threshold tuning} (Nov.~16--23). Thresholds
are tuned on the train split and then frozen.
\item \textbf{Secondary benchmark, BEIR-SciFact} (Nov.~16--23), in parallel.
\item \textbf{Answer generation and QA accuracy} (Nov.~23--30).
\item \textbf{Held-out evaluation on the QASPER dev split and final report}
(Dec.~1--15). The dev split is untouched so far and stays untouched until all
settings are frozen.
\end{itemize}

To absorb the slip we narrow scope in one place: the ablation over reranking
depth becomes a single fixed candidate pool size rather than a sweep. The five
system configurations promised in the proposal are all retained, because the
comparison is the contribution.

\section{Preliminary Results}

\subsection{Experimental Setup}
\label{sec:setup}

\paragraph{Dataset and corpus} QASPER v0.3, train split: 888 papers, 47{,}770
paragraphs, 50{,}189 chunks of 180 words with 30 words of overlap. Both baselines
read the identical \texttt{corpus.jsonl}, verified by SHA-256 against the
recorded build manifest, so the comparison is fair by construction rather than by
convention. Retrieval searches all 888 papers; a question's parent paper is
treated as provenance and never as a candidate filter.

\paragraph{Evaluation subset} QASPER v0.3 is publicly released as train and dev
only, so we use all 2{,}593 train questions for development and reserve the
1{,}005 dev questions as a held-out set to be used after all settings are frozen.
Of the 2{,}593 train questions we score the 1{,}836 (70.8\%) whose every evidence
entry mapped unambiguously to a prepared paragraph. The 757 excluded questions
break down as follows, one mutually exclusive reason each: 271 unanswerable, 7
with annotator disagreement over answerability, 85 with at least one annotation
carrying no evidence, 319 whose evidence points at a figure or table that is not
part of the text corpus, 69 whose evidence is a section heading, and 6 whose
evidence text matches more than one paragraph. Of those 757, 262 do have some
mapped evidence but fail the strict completeness condition; excluding them keeps
the recall denominator honest, at the cost of a smaller sample. Eligible
questions have 1.605 gold paragraphs on average.

\paragraph{Credit granularity} Gold evidence is annotated per paragraph, and our
chunks overlap, so two retrieved chunks can belong to one gold paragraph.
Scoring them as two hits would inflate recall. We therefore fold each retrieved
chunk onto its source paragraph, keep each paragraph's best rank, and then score
at paragraph granularity: $k$ counts distinct paragraphs. We also report the same
cutoffs applied at chunk depth as a cross-check; the two agree at $k = 100$, as
they must.

\paragraph{Metrics} Recall@$k$ for $k \in \{1, 5, 10, 20, 100\}$, averaged over
eligible questions; MRR@10 using the rank of the first gold paragraph; hit
rate@$k$, the fraction of questions with at least one gold paragraph in the top
$k$; and per-query latency. Latency percentiles are recomputed by the evaluator
with a single nearest-rank estimator over one measurement per question, because
each retrieval script records its own percentiles with a different estimator and
those are not mutually comparable.

\paragraph{Systems compared} BM25, our original implementation, $k_1 = 1.5$,
$b = 0.75$, Unicode word tokenizer, average document length 71.9 tokens. Dense
retrieval, \texttt{BAAI/bge-m3}, 1024 dimensions, FAISS \texttt{IndexFlatIP}
over normalized vectors. Both retrieve the top 100 chunks per question with no
reranking and no query rewriting.

\paragraph{Hardware} One Apple Silicon laptop, 10 CPU cores, 32~GB unified
memory, PyTorch 2.15 on the Metal backend. Embedding all 50{,}189 chunks took
37.6~minutes. Dense latency is measured after a warm-up query while BM25 has
none, so the asymmetry favors dense; it applies to one query out of 2{,}593 and
cannot move the reported medians or 95th percentiles.

\subsection{Results}

Table~\ref{tab:results} reports both baselines. Table~\ref{tab:overlap} reports
how often each one finds gold evidence that the other misses.

\begin{table}[tb]
\centering
\small
\caption{Retrieval quality on the QASPER train split, paragraph-level credit over
1{,}836 eligible questions. Both systems use the identical corpus, chunking, and
query set.}
\label{tab:results}
\begin{tabular}{@{}lrr@{}}
\toprule
\textbf{Metric} & \textbf{BM25} & \textbf{Dense (bge-m3)} \\
\midrule
Recall@1       & 0.0262 & \textbf{0.0594} \\
Recall@5       & 0.1018 & \textbf{0.1493} \\
Recall@10      & 0.1419 & \textbf{0.1915} \\
Recall@20      & 0.1918 & \textbf{0.2353} \\
Recall@100     & 0.2924 & \textbf{0.3481} \\
MRR@10         & 0.0742 & \textbf{0.1165} \\
Hit rate@10    & 0.1716 & \textbf{0.2271} \\
Hit rate@100   & 0.3426 & \textbf{0.4020} \\
\midrule
Latency p50 (ms) & \textbf{9.36}  & 23.67 \\
Latency p95 (ms) & \textbf{17.76} & 24.94 \\
\bottomrule
\end{tabular}
\end{table}

\begin{table}[tb]
\centering
\small
\caption{Agreement between the two baselines over the 1{,}836 eligible
questions. Overlap is the mean Jaccard similarity of the retrieved paragraph
sets; the hit columns are the fraction of questions for which that system, and
only that system, retrieves gold evidence.}
\label{tab:overlap}
\begin{tabular}{@{}lrrrrr@{}}
\toprule
 & \textbf{Over-} & \textbf{Both} & \textbf{BM25} & \textbf{Dense} & \textbf{Nei-} \\
\textbf{Depth} & \textbf{lap} & \textbf{hit} & \textbf{only} & \textbf{only} & \textbf{ther} \\
\midrule
$k = 10$  & 0.099 & 0.127 & 0.044 & 0.100 & 0.729 \\
$k = 100$ & 0.098 & 0.279 & 0.063 & 0.123 & 0.535 \\
\bottomrule
\end{tabular}
\end{table}

\subsection{Discussion}

Dense retrieval leads BM25 on every quality metric. At $k = 1$ dense retrieval
is $2.3\times$ better; by $k = 100$ the ratio has fallen to $1.19\times$. MRR@10
sits at $1.57\times$.

A reranker cannot recover evidence absent from the pool it is given, so the
ceiling for our pipeline is the recall of the candidate pool we actually rerank,
which the proposal fixes at 50 chunks --- not Recall@100. At a 50-chunk budget
the paragraph-level hit rate is 0.342 for dense and 0.290 for BM25, against
0.402 and 0.343 at 100 chunks.

The aggregate table does not say whether the two systems find the same evidence,
so we measured that separately (Table~\ref{tab:overlap}). They do not: the mean
Jaccard overlap of their retrieved paragraph sets is 0.099 at $k = 10$ and 0.098
at $k = 100$, so the two candidate pools are close to disjoint. They also fail on
different questions. At $k = 10$, BM25 alone finds gold evidence for 4.4\% of
questions and dense alone for 10.0\%, and at $k = 100$ the union of the two pools
contains gold evidence for 46.5\% of questions against 40.2\% for dense alone.

That union, however, spans 181 paragraphs, and our pipeline reranks 50 chunks.
Splitting a 50-chunk budget evenly between the two retrievers reaches 0.345 ---
indistinguishable from dense alone at the same budget, 0.342, and well short of
the 0.465 the full union reaches. Two consequences follow, and both cut against
the plan in Section~\ref{sec:plan}. Fusion has to earn its gain through ranking
rather than through pool coverage, since at a fixed budget it does not see more
gold than dense alone. And the pool size we fixed at 50 discards roughly a
quarter of the gold that 100 chunks reach, which makes the candidate-pool sweep
we dropped from scope the experiment most worth restoring.

The cost side of the trade is visible too. Dense retrieval is $2.5\times$ slower
at the median. Query encoding accounts for 92\% of that: on 200 queries the
median encode takes 21.8~ms against 1.8~ms for the index search. BM25 does no
neural work at query time. Dense latency is also far more predictable: its p95/p50 ratio
is 1.05 against BM25's 1.90, since BM25's cost scales with how many postings
lists a query touches. This is exactly the regime our adaptive policy targets, and
it tells us the policy should be judged on cross-encoder invocations avoided
rather than on retrieval latency, which is already small relative to reranking.

Absolute values are modest, and the reason is the setup rather than the systems.
Each question is answered against all 888 papers with no same-paper filter, so a
system must first find the right paper among 50{,}189 chunks. QASPER questions
are also written by readers of a specific paper and are frequently generic in
isolation --- ``which baselines did they compare against'' --- which gives a
lexical retriever little to work with across papers. Table~\ref{tab:overlap} is
why we expect fusion to help here.

Two implementation challenges are worth recording. First, the proposal's
intended dense encoder, \texttt{all-MiniLM-L6-v2}, accepts 512 tokens, while our
longest chunk is 1{,}061 tokens; silently truncating a third of a chunk would
have made the dense baseline unfair rather than merely weak. We switched encoders
and made the index builder refuse to run if any chunk would be truncated, so the
failure cannot recur unnoticed. Second, our first evaluator draft credited every
chunk independently, which counted overlapping chunks of one gold paragraph as
separate evidence and overstated recall. The paragraph-level fold in
Section~\ref{sec:setup} fixes this, and a regression test now asserts that two
overlapping chunks of one gold paragraph yield a recall of $1/1$ rather than
$2/1$. A second test asserts that BM25's zero-score padding rows, which appear
when fewer than 100 chunks share any term with the query, cannot earn recall.

\section{Evolving Methodology}

\subsection{Updated Plan}

The architecture is unchanged: retrieve with BM25 and dense retrieval, fuse with
Reciprocal Rank Fusion, estimate the confidence of the fused ranking from
retriever agreement and score separation, and invoke a cross-encoder reranker
only when confidence is low. The contribution remains the adaptive policy that
decides when reranking is necessary. What changed is that every component is now
a committed choice rather than an example, and the evaluation protocol is
specified precisely.

\paragraph{Components reused}
\begin{itemize}
\item \deleted{BM25 implementation e.g.\ \texttt{rank-bm25}} $\rightarrow$ \added{BM25 implemented by us directly, Python standard library
      only, $k_1 = 1.5$, $b = 0.75$}
\item \deleted{pretrained dense model e.g.\ all-MiniLM-L6-v2} $\rightarrow$ \added{dense encoder \texttt{bge-m3}, 1024 dimensions,
      8{,}192-token limit}
\item \deleted{similarity search e.g.\ FAISS} $\rightarrow$ \added{FAISS \texttt{IndexFlatIP} over normalized vectors, exact
      cosine search, no approximation}
\item \deleted{cross-encoder e.g.\
      \texttt{ms-marco-MiniLM-L6-v2}} $\rightarrow$ \added{cross-encoder \texttt{ms-marco-MiniLM-L6-v2}, applied to a fixed
      candidate pool of the top 50 fused chunks}
\item \deleted{LLM for answer generation e.g.\ OpenAI API using a fixed
      ChatGPT model} $\rightarrow$ \added{answer generation by \texttt{qwen3:8b} from the Qwen3
      family, served locally with Ollama, temperature 0, one fixed prompt across
      all configurations}
\item \added{compute platform: one Apple Silicon laptop, 32~GB unified
      memory, PyTorch Metal backend. Embedding the 50{,}189-chunk corpus takes
      37.6~minutes and the cross-encoder is a 22M-parameter model, so GPU
      embedding and reranking do not require PACE. PACE-ICE remains our fallback
      if the secondary benchmark proves larger than this machine can handle.}
\end{itemize}

\paragraph{Chunking and retrieval setup}
\added{Documents are split into chunks of 180 words with 30 words of
overlap, and a chunk never crosses a source paragraph boundary, so each chunk
belongs to exactly one paragraph and carries its identifier. Retrieval searches
the entire split; a question's parent paper is provenance, not a candidate
filter.}

\paragraph{Mapping retrieved chunks to gold evidence}
\added{QASPER annotates gold evidence at paragraph granularity. For
Recall@$K$ and MRR we fold each retrieved chunk onto its source paragraph,
keeping that paragraph's best rank, and score the resulting deduplicated
paragraph ranking against the gold paragraph set, so $K$ counts distinct
paragraphs. Overlapping chunks of one gold paragraph count once. We additionally
report the same cutoffs at chunk depth as a cross-check.}

\paragraph{Evaluation set}
\deleted{For QASPER, we plan to evaluate approximately 1,000 questions
sampled from the designated test set using a fixed random seed.}
\added{QASPER v0.3 is publicly released as train and dev only, with no
evidence-annotated test split. We therefore develop on all 2{,}593 train
questions and hold out all 1{,}005 dev questions for final evaluation after the
adaptive thresholds are frozen. Within a split we score every question whose
evidence mapped completely and unambiguously --- 1{,}836 of 2{,}593 on train,
646 of 1{,}005 on dev --- and report the eligible count together with a
mutually exclusive exclusion reason for every remaining question. No random
sampling is involved, so no seed is needed.}

\paragraph{Metrics}
\deleted{We will report Recall@1, Recall@5, and Recall@10.}
\added{We report Recall@1, @5, @10, @20, and @100, plus MRR@10 and hit
rate@$K$. The deeper cutoffs bound what reranking can reach: a reranker only
reorders the pool it is given, so pool recall at the budget we rerank is the
ceiling, and the gap between it and Recall@10 is what the adaptive policy can
convert. Latency percentiles are recomputed by the
evaluator with one nearest-rank estimator across all systems.}

\subsection{Justification of Changes}

\paragraph{Dense encoder} \texttt{all-MiniLM-L6-v2} accepts 512 tokens. Our
longest chunk is 1{,}061 tokens, so the proposal's encoder would have silently
discarded text from the longest chunks and made the required dense baseline
unfairly weak. \texttt{bge-m3} accepts 8{,}192 tokens, so no chunk is truncated,
and the index builder now refuses to run if any chunk would exceed the limit.

\paragraph{BM25 implemented directly} The project requires non-trivial original
implementation rather than configuration of existing components, and writing
BM25 ourselves also let us make tie-breaking deterministic and bring the index
into the same checksum-verified reproducibility chain as everything else. It
costs little: index construction over 50{,}189 chunks takes 1.9~seconds.

\paragraph{Evaluation set} This change was forced rather than chosen. The
proposal assumed an evidence-annotated QASPER test split; the public v0.3
release has train and dev only. Holding out dev preserves the intent --- tune on
one set, report on another that was never used for tuning.

\paragraph{Eligibility policy} 757 of 2{,}593 train questions cannot support
automatic paragraph-level evaluation, most often because the annotated evidence
is a figure or table that is not in the text corpus at all, or because the
question is marked unanswerable. Scoring those questions would penalize every
system for evidence no text retriever could return. We exclude them, state the
count, and give one reason per excluded question so the denominator is auditable.

\paragraph{Paragraph-level credit} Implementation experience drove this. Our
first evaluator counted chunks independently and overstated recall, because
overlapping chunks of a single gold paragraph were credited as multiple pieces
of evidence. Since QASPER's annotations are paragraph-level, paragraph-level
credit is the measurement the annotations actually support.

\paragraph{Metric depth} Added after seeing the baseline numbers. Pool recall at
the depth we rerank is the ceiling on the whole pipeline, and at the proposal's
cutoffs we could not have measured it.

\subsection{Response to Proposal Feedback}

\paragraph{``Specify the exact chunking strategy and the QASPER retrieval setup
$\ldots$ state how retrieved chunks will be mapped back to gold paragraphs for
Recall@K and MRR.''} Addressed and implemented. Chunks are 180 words with 30
words of overlap and never cross a paragraph boundary; retrieval searches all 888
papers of the split with no same-paper filter. Retrieved chunks are folded onto
their source paragraph, keeping the best rank, and scored against the gold
paragraph set, so overlapping chunks of one gold paragraph count once and $K$
counts distinct paragraphs; chunk-depth figures are reported alongside as a
cross-check. Section~\ref{sec:setup} also states the eligible question count and
every exclusion.

\paragraph{``Commit to concrete choices where the proposal currently says `e.g.'
$\ldots$ name the LLM and how you will access it $\ldots$ and the compute
platform for GPU embedding and reranking.''} Addressed. BM25 is our own
implementation rather than a library; the dense encoder is
\texttt{BAAI/bge-m3} with FAISS \texttt{IndexFlatIP}; the cross-encoder is
\texttt{cross-encoder/\allowbreak ms-marco-MiniLM-L6-v2} over a fixed pool of 50 fused
candidates; answer generation uses \texttt{qwen3:8b} from the Qwen3 family,
served locally with Ollama at temperature 0 with one fixed prompt for all five
configurations, which keeps the generator identical across systems and avoids
an external API dependency. On compute: all work so far ran on one Apple Silicon
laptop using the PyTorch Metal backend, where embedding the full 50{,}189-chunk
corpus took 37.6~minutes, and the cross-encoder has 22M parameters, so GPU
embedding and reranking do not require PACE. We keep PACE-ICE as a documented
fallback for the secondary benchmark.

\section{Potential Risks and Mitigation}

\subsection{Risks Already Addressed}

\paragraph{Unfair comparison between systems} The whole project rests on
comparing five configurations, so any difference in corpus, chunking, or query
set between them would invalidate every conclusion. This is now enforced rather
than promised. Each index records the checksum of the corpus it was built from;
each run records the checksum of the query file and of the index manifest; the
evaluator verifies the entire chain and refuses to score a run produced from a
different query set, and it scores all systems in one invocation so they
necessarily share one eligible set and one latency estimator. An independent
rebuild of the train split reproduced the recorded corpus checksum exactly.

\paragraph{Ambiguity in what counts as a correct retrieval} Resolved by folding
chunks onto paragraphs, fixing a single documented eligibility policy, and
pinning both with regression tests. Two of those tests encode mistakes we
actually made: crediting overlapping chunks separately, and letting BM25's
zero-score padding rows earn recall.

\paragraph{Compute availability for dense retrieval} Initially a schedule risk,
since the dense baseline was the only component needing a GPU. Measured rather
than assumed: 50{,}189 chunks embed in 37.6~minutes on a laptop GPU backend, so
no shared-cluster queue sits on our critical path.

\subsection{Remaining Risks}

\paragraph{Risk 1: the adaptive policy may not beat always reranking, or dense
retrieval alone, by a useful margin.} A cross-encoder has to convert the gap
between Recall@10 and Recall@100 into ranking quality, and the confidence signals
have to identify the queries where reranking pays. Either step can fail. The policy is
also only interesting if it skips a meaningful share of queries while giving up
little quality; a policy that reranks 95\% of queries is not a contribution.

\paragraph{Risk 2: the remaining components may not fit in the time left.} Four
components remain (fusion, reranking with the adaptive policy, the secondary
benchmark, answer generation) against six weeks, and we are already behind on
fusion and the secondary benchmark.

\subsection{Mitigation Plan}

\paragraph{For Risk 1} We report the quality-versus-reranking-rate trade-off as a
curve over threshold values rather than as a single operating point, so the
result is informative even if no threshold dominates always-reranking. The
comparison against always-rerank and against hybrid-without-rerank isolates
where any gain comes from. If the adaptive policy does not beat always-reranking
on quality, the honest and still-publishable result is the trade-off curve plus
an analysis of which confidence signal carried information and which did not;
retriever agreement and score separation are measured separately so we can say
which failed. If the cross-encoder itself turns out not to help, we fall back to
reporting fusion against both baselines and document the negative reranking
result with per-query evidence from the existing per-question output files.

\paragraph{For Risk 2} The components are ordered so that each one is reportable
on its own: fusion first, since it needs no new model and consumes the run files
we already have; then reranking and the adaptive policy, which is the
contribution; then the secondary benchmark; then answer generation, which the
project specification treats as supplementary. If time runs out, we cut from the
back. The secondary benchmark can be reduced to a fixed sampled subset of
SciFact with the sampling procedure reported, and answer generation can be
reported on a sample of questions rather than all of them, since QA accuracy is
a secondary metric and the generator is held constant across systems either way.
Our evaluation harness already accepts any number of run files in one
invocation, so adding each new configuration to the results table costs one
command rather than new analysis code.

\end{document}
